\section{ماژول dobi\_svd (Dobi-SVD Compressor)}
\label{sec:module_dobi_svd}

\subsection{معرفی و انگیزه}
این ماژول در چارچوب پژوهش حاضر پیاده‌سازی و ارزیابی شده و فاقد مقاله مستقل جداگانه در ادبیات تحقیق است.

در روش‌های تقریب رتبه پایین مرسوم مبتنی بر \lr{SVD}، رتبه تقریب $r$ به صورت یک ابرپارامتر گسسته و ثابت در تمامی لایه‌ها تعیین می‌شود. این رویکرد غیرانعطاف‌پذیر به دلیل ناهمگونی ساختاری لایه‌ها با شکست یا افت بهره‌وری مواجه می‌گردد؛ چراکه برخی لایه‌ها به ظرفیت بازنمایی با رتبه بالا وابسته هستند در حالی که برخی دیگر افزونگی بسیار شدیدی دارند و می‌توان آن‌ها را به رتبه‌های بسیار پایین تقلیل داد. علاوه بر آن، هرس ناگهانی مؤلفه‌های تکین پس از فاکتورگیری امکان تطبیق پیوسته وزن‌ها با توزیع داده‌های هدف را از بین می‌برد.

ماژول فشرده‌سازی \lr{Dobi-SVD Compressor} به منظور حل ساختاری این مسئله از طریق معرفی لایه‌های خطی مشتق‌پذیر مبتنی بر دریچه‌بندی طیفی پیوسته (\lr{Differentiable SVD with Learnable Spectral Gating}) ابداع شده است. در این معماری، فاکتورهای کم‌مرتبه با بردار دروازه طیفی پیوسته و آموزش‌پذیر $\boldsymbol{\alpha}$ ترکیب می‌شوند که به مدل اجازه می‌دهد در طول فاز بهینه‌سازی گرادیانی، سهم هر مؤلفه طیفی را به صورت مشتق‌پذیر تعدیل کرده یا به سمت صفر میل دهد. این امر به هرس خودکار و پویای رتبه (\lr{Differentiable Rank Pruning}) منجر گردیده و بالاترین پایداری را در فاز بازیابی ایجاد می‌کند.

\subsection{مبانی نظری و فرمول‌بندی ریاضی}
فرض کنید ماتریس تبدیل خطی $\mathbf{W} \in \mathbb{R}^{d_{\text{out}} \times d_{\text{in}}}$ با رتبه حداکثری اولیه $r \le \min(d_{\text{out}}, d_{\text{in}})$ تقریب زده شود. لایه خطی اختصاصی \lr{DobiSVDLinearLayer} ماتریس وزن موثر بازسازی‌شده را به صورت حاصل‌ضرب فاکتورهای متعامد مقیاس‌گذاری‌شده با بردار دروازه طیفی $\boldsymbol{\alpha} \in \mathbb{R}^r$ تعریف می‌کند:
\begin{equation}
\mathbf{W}_{\text{eff}}(\boldsymbol{\alpha}) = \mathbf{U} \cdot \left( \operatorname{diag}(\sigma(\boldsymbol{\alpha})) \, \mathbf{V} \right)
\end{equation}
که در آن:
\begin{itemize}
    \item ماتریس $\mathbf{U} \in \mathbb{R}^{d_{\text{out}} \times r}$ و ماتریس $\mathbf{V} \in \mathbb{R}^{r \times d_{\text{in}}}$ پارامترهای آموزش‌پذیر وزن هستند که در آغاز با فاکتورهای متقارن حاصل از تجزیه اولیه مقدار منفرد مقداردهی می‌شوند:
    \begin{equation}
    \mathbf{W} = \mathbf{P} \mathbf{\Sigma} \mathbf{Q}^\top, \quad \mathbf{U}_{\text{init}} = \mathbf{P}_{:, :r} \mathbf{\Sigma}_r^{1/2}, \quad \mathbf{V}_{\text{init}} = \mathbf{\Sigma}_r^{1/2} \mathbf{Q}_{:, :r}^\top
    \end{equation}
    \item تابع $\sigma(\cdot)$ نشان‌دهنده تابع غیرخطی سیگموئید است:
    \begin{equation}
    \sigma(\alpha_k) = \frac{1}{1 + e^{-\alpha_k}}, \quad k \in \{1, \dots, r\}
    \end{equation}
    بردار دروازه $\boldsymbol{\alpha}$ به صورت پیش‌فرض با مقدار اولیه مثبت $\alpha_k^{(0)} = 2.5$ مقداردهی می‌شود تا در آغاز فرآیند، ضریب دروازه مقداری نزدیک به یک داشته باشد ($\sigma(2.5) \approx 0.924$) و از شوک اولیه‌ی هرس به مدل جلوگیری شود.
\end{itemize}

در گذر پیشرو، خروجی به شکل بهینه‌شده و بدون تشکیل صریح ماتریس کامل وزن محاسبه می‌گردد:
\begin{equation}
\mathbf{y} = \left( \mathbf{x} \left( \operatorname{diag}(\sigma(\boldsymbol{\alpha})) \mathbf{V} \right)^\top \right) \mathbf{U}^\top + \mathbf{b} = \left( (\mathbf{x} \mathbf{V}^\top) \odot \sigma(\boldsymbol{\alpha})^\top \right) \mathbf{U}^\top + \mathbf{b}
\end{equation}

این رابطه نشان می‌دهد که عملگر دروازه‌بندی با هزینه محاسباتی ناچیز ضرب عنصری $\mathcal{O}(r)$ بر روی فعال‌سازی‌های پنهان میانی اعمال می‌شود و مشتقات نسبت به تمامی پارامترها به صورت کاملاً تحلیلی و ملایم منتشر می‌گردند:
\begin{equation}
\frac{\partial \mathcal{L}}{\partial \alpha_k} = \sigma(\alpha_k)(1 - \sigma(\alpha_k)) \cdot \left( \frac{\partial \mathcal{L}}{\partial (\mathbf{W}_{\text{eff}})_{:, :}} \right) : (\mathbf{U}_{:, k} \mathbf{V}_{k, :})
\end{equation}

\subsection{معماری و پیاده‌سازی}
ماژول \lr{Dobi-SVD} در فایل \lr{Codes/framework/compression\_engine/dobi\_svd.py} پیاده‌سازی شده و تحت شناسه‌های \lr{dobi} و \lr{dobi\_svd} در رجیستری فشرده‌سازی ثبت گردیده است. ارکان ساختاری آن شامل موارد زیر است:
\begin{itemize}
    \item \textbf{ماژول خطی \lr{DobiSVDLinearLayer}:} پیاده‌سازی کاملاً ماژولار و مبتنی بر \lr{nn.Module} که جایگزین مستقیم لایه‌های \lr{nn.Linear} و لایه‌های \lr{Conv1D} در مدل‌های زبانی می‌شود.
    \item \textbf{انطباق پویا با طول ابعاد پنهان (\lr{Dimension Padding}):} لایه مجهز به منطق پدینگ پویا است تا در صورت عدم تطابق طول بردار ورودی، داده‌ها با صفر گسترش یافته و از خطای ابعادی جلوگیری شود.
    \item \textbf{فیلتر هوشمند لایه‌ها و پیش‌گیری از شکست:} نادیده گرفتن لایه‌های غیرخطی، لایه‌های هنجارسازی (\lr{LayerNorm}) و پروجکشن‌های طبقه‌بند خروجی، تضمین‌کننده حفظ توانایی‌های زبانی مدل است.
    \item \textbf{پشتیبانی از ارتقای رتبه و خطای بازسازی فروبنیوس:} در صورت فراتر رفتن خطای بازسازی از حد مجاز (\lr{max\_error=0.35})، فرآیند ارتقای رتبه تا دو مرتبه انجام می‌پذیرد.
\end{itemize}

\subsection{نتایج تجربی}
ارزیابی ماژول \lr{dobi\_svd} بر اساس داده‌های تاییدشده در \lr{thesis\_final/master\_verified\_results.json} روی مدل‌های پایه به شرح مندرج در جدول \ref{tab:benchmark_dobi_svd} به ثبت رسیده است:

\begin{table}[H]
\centering
\small
\caption{نتایج بنچ‌مارک سخت‌افزاری و زبانی فشرده‌ساز \lr{dobi\_svd} روی مدل‌های پایه}
\label{tab:benchmark_dobi_svd}
\begin{tabular}{lccccc}
\toprule
\textbf{مدل پایه \lr{(Base Model)}} & \textbf{دقت \lr{(Accuracy)}} & \textbf{شاخص \lr{F1}} & \textbf{حافظه \lr{(Peak VRAM)}} & \textbf{توان \lr{(Throughput)}} & \textbf{وضعیت} \\
\midrule
\lr{DistilBERT} & \lr{90.48\%} & \lr{0.9079} & \lr{310.69 MB} & \lr{52,426.55 tok/s} & \lr{PASS} \\
\lr{RoBERTa} & \lr{74.54\%} & \lr{0.7707} & \lr{530.19 MB} & \lr{27,881.66 tok/s} & \lr{PASS} \\
\lr{GPT-2} & \lr{58.14\%} & \lr{0.6256} & \lr{635.79 MB} & \lr{27,358.79 tok/s} & \lr{PASS} \\
\lr{Qwen} & \lr{89.45\%} & \lr{0.8950} & \lr{1850.64 MB} & \lr{4,747.76 tok/s} & \lr{PASS} \\
\bottomrule
\end{tabular}
\end{table}

داده‌های تجربی جدول \ref{tab:benchmark_dobi_svd} حاکی از عملکرد قابل توجه و همه‌جانبه ماژول \lr{dobi\_svd} است که در هر چهار مدل پایه موفق به کسب نشان رسمی \lr{PASS} گردیده است. به طور خاص:
\begin{itemize}
    \item در مدل \lr{DistilBERT}، دقت رکوردشکن \lr{90.48\%} و شاخص اف‌وان \lr{0.9079} همراه با توان پردازشی بالای \lr{52,426.55 tok/s} و مصرف حافظه ناچیز \lr{310.69 MB} ثبت شده است.
    \item در مدل زبانی نسل جدید \lr{Qwen}، این الگوریتم به دقت خیره‌کننده \lr{89.45\%} و شاخص اف‌وان \lr{0.8950} با توان \lr{4,747.76 tok/s} دست یافته که نشان‌دهنده اثربخشی بالا دروازه‌بندی طیفی در مدل‌های بسیار بزرگ با ساختار علّی است.
    \item در مدل \lr{RoBERTa} و \lr{GPT-2} نیز با وجود پیچیدگی‌های معماری، وضعیت \lr{PASS} با دقت‌های به ترتیب \lr{74.54\%} و \lr{58.14\%} به دست آمده است.
\end{itemize}

\subsection{محدودیت‌ها و تحلیل شکست}
با وجود نتایج بسیار قوی، تحلیل مکانیزم‌های عملکردی ماژول نقاط حساس و محدودیت‌های زیر را برجسته می‌سازد:
\begin{itemize}
    \item \textbf{نیاز به گام‌های بهینه‌سازی گرادیانی جهت فعال‌سازی گیت‌ها:} در صورتی که پس از اعمال تقریب \lr{Dobi-SVD}، فرصت کافی برای به‌روزرسانی پارامترهای $\boldsymbol{\alpha}$ از طریق داده‌های بازیابی فراهم نشود، مقادیر اولیه دروازه ($\sim 0.92$) ممکن است انحراف مقیاسی کوچکی در وزن‌ها پدید آورند که علت افت نسبی دقت در \lr{GPT-2} (\lr{58.14\%}) به همین مسئله مربوط است.
    \item \textbf{حساسیت مقادیر گرادیان دروازه‌ها در رتبه‌های کوچک:} مقادیر منفرد انتهایی با مقیاس‌های بسیار کوچک ممکن است گرادیان‌های ناپایداری تولید کنند، از این رو لزوم فیلترسازی و منظم‌سازی دقیق رتبه‌های نزدیک به صفر الزامی است.
    \item \textbf{سربار حافظه در مدل‌های دیکودر:} تخصیص پارامترهای اختصاصی دروازه و عملیات گسترش عنصری در لایه‌های تو در تو تا حدودی اوج مصرف حافظه را در مقایسه با \lr{Direct SVD} افزایش می‌دهد.
\end{itemize}
